\hypertarget{appendix-c99977d88b33}{}
\section*{Unsigned modular arithmetic}
\section{Unsigned modular arithmetic}\label{unsigned-modular-arithmetic}

\textbf{Curriculum depth:} Developed · \textbf{First seen:} C001 · \textbf{Latest development:} C002

\subsection{Definition and mechanism}\label{definition-and-mechanism}

Conversion and arithmetic in an unsigned \(N\)-bit domain produce residues modulo \(2^N\); intentional wrap is valid, but wrapped bounds expressions can become security bugs.

Modulo m retains the remainder class after division by m. An unsigned N-bit domain uses m=2ᴺ: 250+10 becomes 4 when stored in eight bits, and 0−1 becomes 255. Narrow operands may promote first, so identify the operation type separately from the destination type.

\subsection{Prerequisites}\label{prerequisites}

\begin{itemize}
\tightlist
\item
  \hyperref[appendix-83c3f28f7bda]{Positional notation} --- A base is the number of available digits.
\item
  \hyperref[appendix-d67a19507dc5]{Machine word} --- A processor's natural integer processing width is often called its word width, but the term is architecture dependent.
\end{itemize}

\subsection{Used by}\label{used-by}

\begin{itemize}
\item
  \hyperref[appendix-53a8cfb25f8d]{Offset-binary coding} --- uses this mechanism to define its representation, error or validation contract.
\item
  \hyperref[appendix-652d81568de9]{Checked binary parsing} --- uses this mechanism to define its representation, error or validation contract.
\item
  \hyperref[appendix-4fbc03d20829]{Two\textquotesingle s complement} --- uses unsigned modular arithmetic as part of its representation or reasoning.
\item
  \hyperref[appendix-2ed044358b82]{Arithmetic policies} --- uses unsigned modular arithmetic as part of its representation or reasoning.
\end{itemize}

\subsection{Recall and next depth}\label{recall-and-next-depth}

Explain unsigned modular arithmetic without the label, work the example above, then state the assumption whose removal would change the result.

\textbf{Pending:} modular algorithms, counters, hashes, and sequence-number comparisons.

\subsection{Source and navigation}\label{source-and-navigation}

\href{Continuous\%20Notes/Computing\%20Foundations\%20-\%20From\%20Bits\%20to\%20C++\%20and\%20Systems\%20APIs\#From\%20a\%20signal\%20to\%20an\%20interpretation}{C001 teaching chapter} · \href{Continuous\%20Notes/Computing\%20Foundations\%20-\%20From\%20Bits\%20to\%20C++\%20and\%20Systems\%20APIs\#Arithmetic\%20becomes\%20a\%20policy}{C002 development} · \hyperref[appendix-fd090d982a15]{Part definitions and mechanisms} · \href{}{Dependency Map (02-Foundations-II, appendix-4ad47bdd269a)}
